\chapter*{Conclusion générale}
\markboth{Conclusion générale}{}
\addcontentsline{toc}{chapter}{Conclusion générale}

\section*{Bilan du projet}

Ce projet de fin d'études a abouti à la réalisation complète de
\textbf{LeadGen Francophone 360+}, une plateforme de génération automatisée de
leads B2B IT destinée à Solvinya Group. Tous les objectifs initiaux ont été
atteints~:

\begin{itemize}
  \item \textbf{Collecte multi-source~:} 4 collecteurs actifs (LinkedIn Jobs,
        LinkedIn B2B, BOAMP, TED EU) couvrant 5 pays francophones (France,
        Belgique, Luxembourg, Suisse, Tunisie) et produisant \textbf{815 prospects}
        lors de la première exécution complète~;
  \item \textbf{Scoring déterministe~:} formule pondérée 4 dimensions (pays,
        technologie, secteur, boost) produisant des scores 0--100 cohérents
        avec les priorités métier de Solvinya~;
  \item \textbf{Module NLP~:} classifieur TF-IDF + LinearSVC atteignant
        \textbf{75\%} de précision en cross-validation 5-fold sur 10 secteurs,
        surpassant les embeddings sentence-transformers sur ce corpus technique~;
  \item \textbf{Générateur LLM~:} messages LinkedIn et emails personnalisés
        générés en 1 à 3 secondes via Groq, adaptés au type de signal détecté~;
  \item \textbf{Automatisation~:} 4 workflows n8n opérationnels, collecte
        planifiée toutes les 12h, déclenchement manuel par webhook~;
  \item \textbf{API REST~:} 8 endpoints FastAPI documentés, authentifiés et
        testés~;
  \item \textbf{Infrastructure~:} déploiement Docker Compose en une commande,
        dashboard Metabase connecté.
\end{itemize}

\section*{Difficultés rencontrées}

Plusieurs difficultés techniques ont été surmontées au cours du projet~:

\begin{description}
  \item[API TED EU] L'API ne fonctionne qu'en POST (le GET retourne 405). La
    syntaxe \emph{expert query} avec des noms de champs très spécifiques
    (\texttt{notice-title}, \texttt{buyer-country}) a nécessité une exploration
    de la documentation officielle. De plus, chaque requête ne peut filtrer
    qu'un seul pays, imposant une boucle Python.

  \item[LinkedIn CONTENT filter] L'API non-officielle \texttt{linkedin-api} ne
    retourne aucun résultat avec le filtre \texttt{CONTENT}. La stratégie a été
    pivotée vers la recherche d'entreprises acheteuses suivie du scan de leurs
    publications.

  \item[Malt.fr] La protection Cloudflare bloque toutes les requêtes httpx et
    même Playwright standard. Ce collecteur reste inactif dans la version
    actuelle.

  \item[Déduplication asyncpg] Un bug de paramétrage SQL (mélange de valeurs
    littérales et de paramètres \texttt{\$N}) a causé des erreurs de comptage de
    paramètres, corrigées en rendant tous les paramètres dynamiques.
\end{description}

\section*{Perspectives d'évolution}

\begin{enumerate}
  \item \textbf{Malt.fr~:} implémentation avec Playwright + stealth plugin pour
        contourner Cloudflare~;
  \item \textbf{Enrichissement contacts~:} intégration Hunter.io pour trouver
        l'email du DSI ou DRH automatiquement~;
  \item \textbf{Qualification par LLM~:} utiliser un LLM pour pré-qualifier les
        prospects (est-ce réellement un besoin Solvinya~?) avant le scoring~;
  \item \textbf{CRM intégration~:} connecter la plateforme à HubSpot ou Pipedrive
        via API pour synchroniser les prospects qualifiés directement~;
  \item \textbf{Monitoring~:} ajouter Prometheus + Grafana pour surveiller le
        volume de collecte, les taux d'échec et les latences API~;
  \item \textbf{Multi-tenant~:} étendre la plateforme pour servir plusieurs ESN
        avec des paramètres de scoring configurables par client.
\end{enumerate}

\section*{Conclusion personnelle}

Ce projet m'a permis d'appliquer concrètement des compétences transversales~:
développement Python asynchrone (asyncio, asyncpg, FastAPI), traitement
automatique du langage naturel, intégration d'APIs complexes, et DevOps
(Docker, n8n). Au-delà des aspects techniques, j'ai appris à pivoter rapidement
face aux limitations des APIs tierces et à concevoir un système résilient dont
chaque composant peut fonctionner indépendamment.

La plateforme est opérationnelle, documentée et deployable, et constitue une
base solide pour l'automatisation de la prospection commerciale de Solvinya Group.
